\begin{figure}[t]
\centering
\resizebox{\columnwidth}{!}{%
\begin{tikzpicture}
\begin{axis}[
  width=8.4cm, height=5.7cm,
  xlabel={constellation size (satellites)}, ylabel={MAE (lower is better)},
  symbolic x coords={66, 264, 1584}, xtick=data,
  ymin=0.04, ymax=0.23,
  legend pos=north west, legend cell align=left,
  grid=major, grid style={gray!18},
  tick label style={font=\footnotesize}, label style={font=\small},
  legend style={font=\scriptsize, fill=white, fill opacity=0.9, draw=gray!40},
  axis line style={gray!50},
]
\addplot[gray!55, very thick, mark=triangle*, mark size=2pt] coordinates {
  (66,0.102) (264,0.163) (1584,0.208)};
\addlegendentry{GCN (local)}
\addplot[steel, very thick, mark=*, mark size=2pt] coordinates {
  (66,0.064) (264,0.069) (1584,0.168)};
\addlegendentry{Heat (classical global)}
\addplot[amber!80!black, very thick, mark=square*, mark size=2pt] coordinates {
  (66,0.085) (264,0.071) (1584,0.148)};
\addlegendentry{Quantum-inspired walk}
\node[font=\scriptsize\bfseries, text=amber!75!black, align=center]
  at (axis cs:1584,0.118) {QI best};
\end{axis}
\end{tikzpicture}}
\caption{Exp~E, the quantum-inspired edge emerges at scale. Operator MAE versus
constellation size at a fixed training budget. The quantum-inspired walk (amber) starts
worse than the classical heat kernel (blue) at $66$ satellites, ties at $264$, and
overtakes it at $1584$, where geodesics are longest and the ballistic ($\Var\!\sim\!t^2$)
reach of the walk pays off; the local GCN trails at every scale.}
\label{fig:scale}
\end{figure}
